\section{程序设计与实现}

程序入口为 \codepath{task2/src/stereo_depth.py}，流程包括输入检查、同步预处理、
StereoSGBM 匹配、有效视差筛选、相对深度计算和结果保存。

\subsection{图像读取与预处理}

程序使用 \texttt{cv2.imread} 读取左右彩色图像，并检查两幅图像均成功读取且尺寸一致。
随后将图像转为灰度图用于匹配。默认配置对左右彩色图和灰度图同步执行一次
\texttt{cv2.pyrDown}，得到 $641\times555$ 的处理图像。

\subsection{StereoSGBM 视差计算}

实际运行参数来自 \codepath{task2/results/params.json}，列于表~\ref{tab:sgbm-params}。

\begin{table}[H]
  \centering
  \caption{本次运行使用的 StereoSGBM 参数}
  \label{tab:sgbm-params}
  \begin{tabular}{lr}
    \toprule
    参数 & 实际值 \\
    \midrule
    \texttt{window\_size} & $3$ \\
    \texttt{minDisparity} & $16$ \\
    \texttt{numDisparities} & $96$ \\
    \texttt{blockSize} & $5$ \\
    $P_1$ & $72$ \\
    $P_2$ & $288$ \\
    \texttt{disp12MaxDiff} & $1$ \\
    \texttt{uniquenessRatio} & $10$ \\
    \texttt{speckleWindowSize} & $100$ \\
    \texttt{speckleRange} & $32$ \\
    \bottomrule
  \end{tabular}
\end{table}

OpenCV 匹配器输出为带固定缩放的视差，程序除以 $16$ 转换为像素单位的浮点数组。
本次有效条件为有限值且 $d>15$；无效位置保留为 `NaN`，不参与深度统计。核心逻辑为：

\begin{lstlisting}[language=Python]
disparity = matcher.compute(left_match, right_match).astype(np.float32) / 16.0
valid = np.isfinite(disparity) & (disparity > 15)
disparity_raw = np.where(valid, disparity, np.nan)
depth_proxy = np.where(valid, 1.0 / disparity_raw, np.nan)
\end{lstlisting}

视差图、有效掩码和相对深度图的 PNG 仅用于显示，显示过程使用有效像素归一化，
不能替代原始浮点数组。

\subsection{输出文件}

程序实际保存了以下结果：

\begin{itemize}
  \item \codepath{left.png}、\codepath{right.png}：实际读入并经同步降采样的左右图；
  \item \codepath{disparity_raw.npy}：像素单位的浮点视差，无效区域为 `NaN`；
  \item \codepath{disparity.png}：归一化视差可视化；
  \item \codepath{valid_mask.png}：有效视差区域掩码；
  \item \codepath{depth_proxy.npy} 和 \codepath{depth_proxy.png}：任意尺度相对深度；
  \item \codepath{params.json}：输入、预处理、匹配参数、运行时版本和输出摘要。
\end{itemize}
